\section{Data and research framework}
\label{sec:data}

\subsection{Prices and sources}

The monthly dataset covers January 2014-January 2026, including benchmark,
membership, sector and corporate-action records. We label each month's latest
observation by calendar month-end and prioritize Reuters, unadjusted Yahoo
\texttt{Close}, then a fixed, hash-verified WIKI extract. Missing latest Reuters
closes never revive stale quotes. Duplicate Reuters date/security rows merge
only if nonmissing prices and volumes agree; conflicts fail preparation.
We neither interpolate prices nor replace historical constituents with today's
universe.

\begin{table}[htbp]
\centering
\begin{threeparttable}
\small
\begin{tabularx}{\linewidth}{@{}>{\raggedright\arraybackslash}p{0.24\linewidth}X@{}}
\toprule
\textbf{Source or evidence} & \textbf{Role and interpretation} \\
\midrule
Supplied Reuters & Primary monthly backtest prices and volumes; reviewed
fallbacks fill identified gaps. \\
Yahoo Finance & S\&P 500 Total Return benchmark, adjusted daily competition
replay prices, and specific backtest or shares requirements. \\
Historical membership & Effective-dated S\&P 500 baskets from the
\texttt{fja05680/sp500} history, checked against changes and intervals. \\
Dated Wikipedia revisions & Historical GICS sector evidence, supplemented by
reviewed S\&P addition notices only where a new member is absent. \\
Issuer / SEC evidence & Security events, transaction terms, issuer
continuity and reviewed outstanding-share observations. \\
\bottomrule
\end{tabularx}
\caption[Data sources and analytical roles]{Data sources and their analytical roles.}
\label{tab:data-sources}
\end{threeparttable}
\end{table}

The backtest uses mixed-source \emph{price returns} with reviewed corporate
actions but omits ordinary dividends. The competition replay uses adjusted
Yahoo prices and reviewed monthly reconstructions, including dividends once
in closes or event-return factors. This difference limits comparison with the
total-return benchmark. Each result records its price convention and
simulation fingerprint.

\subsection{Identities and historical coverage}

A canonical \texttt{Asset\_ID} distinguishes ticker changes from conversions
into different securities. Mappings are start-inclusive and end-exclusive;
overlaps, orphan references and unapproved aliases fail validation. Historical
membership baskets, changes and intervals must agree at each decision date.
The membership history was reconstructed after the competition.

Sectors use the final Wikipedia revision before 00:00 UTC on each required
date. Reviewed S\&P notices fill missing new members from their effective dates
without overwriting classifications; conflicts fail. Assignments retain GICS
codes, source-era labels and references. This reconstruction is not a licensed
GICS history.

Structured corporate-action records govern signed entitlements, fractional
shares, restricted proceeds and event-date financing. Mandatory delivery
incurs no trade fee; a required sale of an off-universe security does.
Attribution capitalization uses prices, capitalization factors and reviewed
shares, preferring SEC/issuer counts, then event estimates, then consistent
Yahoo responses. Missing counts stop preparation; later publications may
establish earlier share counts for retrospective attribution.

Manifests and audits record hashes and transformations; raw provider data
remain local. Backtest eligibility requires a positive holding-period endpoint
price or approved event valuation. Although signals use only lagged history,
this availability filter is retrospective and can favor simulated
feasibility.
